\section{声子介导的有效相互作用}

Gross 第 13 章用二阶微扰，从电--声哈密顿量逐步导出电子之间的有效相互作用；在费米面附近可为吸引，从而触发 Cooper 不稳定性。

\subsection{模型}

取
\begin{equation}
  H
  =H_{\mathrm{el}}+H_{\mathrm{ph}}+H_{\mathrm{ep}},
\end{equation}
其中 $H_{\mathrm{ep}}=\sum_{\mathbf{k}\mathbf{q}\sigma}M_{\mathbf{q}}
c_{\mathbf{k}+\mathbf{q}\sigma}^\dagger c_{\mathbf{k}\sigma}
(a_{\mathbf{q}}+a_{-\mathbf{q}}^\dagger)$（Fr\"ohlich；$|M_{\mathbf{q}}|^2$ 即前文 $|V_{\mathbf{q}}|^2$）。无扰真空取为无声子、电子占据给定组态。

\subsection{单电子自能（二阶）}

对状态 $|\mathbf{k}\sigma\rangle=c_{\mathbf{k}\sigma}^\dagger|\mathrm{vac}\rangle$，中间态为发射一个 $(\mathbf{q})$ 声子后电子到 $\mathbf{k}+\mathbf{q}$：
\begin{equation}
  |\mathrm{int}\rangle
  =a_{\mathbf{q}}^\dagger c_{\mathbf{k}+\mathbf{q}\sigma}^\dagger|\mathrm{vac}\rangle,
  \qquad
  E_{\mathrm{int}}
  =\xi_{\mathbf{k}+\mathbf{q}}+\omega_{\mathbf{q}}.
\end{equation}
二阶能量移动
\begin{equation}
  (\Delta E)^{(2)}_{\mathbf{k}}
  =\sum_{\mathbf{q}}
  \frac{|M_{\mathbf{q}}|^2}
  {\xi_{\mathbf{k}}-(\xi_{\mathbf{k}+\mathbf{q}}+\omega_{\mathbf{q}})}.
\label{eq:polaron_shift}
\end{equation}
这重整化单电子色散（极化子移动）。吸收声子的过程对 $T=0$ 真空不贡献（无热声子）。

\subsection{两电子：交换声子}

对 $|\mathbf{k}\sigma,\mathbf{k}'\sigma'\rangle$，除各自 \eqref{eq:polaron_shift} 外，还有交叉中间态：一个电子发射 $\mathbf{q}$、另一个吸收。矩阵元非零要求动量匹配。该交叉项可改写为对有效势的一阶贡献：
\begin{equation}
\begin{aligned}
  &\langle\mathbf{k}+\mathbf{q},\mathbf{k}'-\mathbf{q}|
  V_{\mathrm{eff}}|
  \mathbf{k},\mathbf{k}'\rangle\\
  &\quad=
  |M_{\mathbf{q}}|^2
  \Biggl[
  \frac{1}{\xi_{\mathbf{k}}-(\xi_{\mathbf{k}+\mathbf{q}}+\omega_{\mathbf{q}})}
  +\frac{1}{\xi_{\mathbf{k}'}-(\xi_{\mathbf{k}'-\mathbf{q}}+\omega_{\mathbf{q}})}
  \Biggr].
\end{aligned}
\end{equation}
若能量转移记为 $\omega$（把时间依赖/推迟结构恢复），并取两电子都在费米面附近，则合并为
\begin{equation}
  V_{\mathrm{eff}}(\mathbf{q},\omega)
  =\frac{2|M_{\mathbf{q}}|^2\omega_{\mathbf{q}}}
  {\omega^2-\omega_{\mathbf{q}}^2}.
\label{eq:Veff_ph}
\end{equation}
（与用声子传播子 $\mathcal{D}(\mathbf{q},\omega)$ 写出的 $-|M|^2\mathcal{D}$ 同构。）当 $|\omega|<\omega_{\mathbf{q}}$ 时 $V_{\mathrm{eff}}<0$：\textbf{延迟吸引}。恰在 $\varepsilon_F$ 上、$\omega\to0$ 时
\begin{equation}
  V_{\mathrm{eff}}(\mathbf{q},0)
  =-\frac{2|M_{\mathbf{q}}|^2}{\omega_{\mathbf{q}}}.
\end{equation}

\subsection{与库仑的竞争}

总有效相互作用 $\approx v_{\mathbf{q}}/\varepsilon(\mathbf{q},\omega)+V_{\mathrm{eff}}$。静态库仑被屏蔽后仍偏排斥，但在 Debye 窗口 $|\omega|\lesssim\omega_D$ 内声子吸引可占优，从而 $N(0)|V|>0$ 的 BCS 参数。Eliashberg 保留完整 $\omega$ 依赖而不用瞬时 $V$。

\subsection{吸引不等于真空束缚}

三维势阱中，吸引势须超过临界强度才形成双体束缚态。费米海上的 Cooper 问题不同：费米海的 Pauli 阻塞使散射振幅出现 $\ln$ 发散，\textbf{任意弱}吸引都产生不稳定性。因此“声子给出 $V_{\mathrm{eff}}<0$”只是必要条件；超导性来自 Cooper 对数，而非真空中的双电子束缚。
